\documentclass[11pt]{article}
\usepackage[T1]{fontenc}
\usepackage{lmodern,amsmath,amssymb,amsthm,mathtools}
\usepackage[margin=1in]{geometry}
\usepackage{microtype,enumitem,xcolor}
\usepackage[numbers,sort&compress]{natbib}
\usepackage{hyperref}
\hypersetup{colorlinks=true,linkcolor=blue!45!black,citecolor=blue!45!black,
urlcolor=blue!45!black,pdftitle={Extremizers for the independent Berry-Esseen inequality},pdfauthor={Logan Bell}}
\setlength{\emergencystretch}{2em}
\numberwithin{equation}{section}
\newtheorem{theorem}{Theorem}[section]
\newtheorem{proposition}[theorem]{Proposition}
\newtheorem{lemma}[theorem]{Lemma}
\newtheorem{corollary}[theorem]{Corollary}
\theoremstyle{definition}
\newtheorem{definition}[theorem]{Definition}
\theoremstyle{remark}
\newtheorem{remark}[theorem]{Remark}
\DeclareMathOperator{\Var}{Var}
\DeclareMathOperator{\diam}{diam}
\DeclareMathOperator{\supp}{supp}
\DeclareMathOperator{\Law}{Law}
\DeclareMathOperator{\Bern}{Bern}
\DeclareMathOperator{\Bin}{Bin}
\DeclareMathOperator{\sinc}{sinc}
\newcommand{\E}{\mathbb E}
\newcommand{\Pp}{\mathbb P}
\newcommand{\R}{\mathbb R}
\newcommand{\Z}{\mathbb Z}
\newcommand{\ind}{\mathbf1}
\newcommand{\ce}{C_{\mathrm E}}
\newcommand{\ci}{C_{\mathrm{ind}}}
\newcommand{\phz}{\phi(0)}
\newcommand{\eps}{\varepsilon}
\newcommand{\TV}{\mathrm{TV}}
\title{Extremizers for the independent Berry--Esseen inequality}
\author{Logan Bell}
\date{11 September 2026}
\begin{document}
\maketitle
\begin{abstract}
Suppose the optimal Berry--Esseen constant for arbitrary independent
summands exceeds Esseen's constant. We prove that the supremum is attained
in a class of countable independent convolutions with an infinitely
divisible remainder, keeping the sum of third absolute moments as part of
the representation. A maximizing representation can be chosen with
three-point summands. Its Gaussian and L\'evy components vanish, and the
supports of its independent summands shrink to zero. If infinitely many
summands are nonconstant, their sum is nonatomic. We also prove a rank
condition on the finite support configurations that can complete to a
maximizing threshold. This condition excludes representations by infinitely
many separated finite blocks, without restrictions on their probabilities
or cardinalities. The proofs use exact variations of the original laws;
CDF discontinuities are retained throughout. The sharp inequality for
all independent arrays remains open.
\end{abstract}

\section{Statements and notation}\label{sec:statements}
For a centered random variable $S$ of variance one, set
\[
 \Delta(S)=\sup_{x\in\R}|\Pp(S\le x)-\Phi(x)|.
\]
The standard normal density and distribution function are denoted by
$\phi$ and $\Phi$. Define
\[
 \ci=\sup\frac{\Delta(\sum_{i=1}^nX_i)}
                        {\sum_{i=1}^n\E|X_i|^3},
 \qquad
 \ce=\frac{3+\sqrt{10}}{6\sqrt{2\pi}},
\]
where the supremum is over finite independent centered arrays with total
variance one and finite third absolute moments. The classical
Berry--Esseen theorem gives $\ci<\infty$, and Esseen's two-point
asymptotics give $\ci\ge\ce$ \cite{Esseen1956,Shevtsova2013}.

We use the small-Lyapunov theorem proved in the companion manuscript
\cite[Theorem~1.1]{BellSmall2026}: there is $\ell_0>0$ such that
\begin{equation}\label{eq:small-input}
 \sum_i\E|X_i|^3\le\ell_0
 \quad\Longrightarrow\quad
 \Delta(\textstyle\sum_iX_i)\le\ce\sum_i\E|X_i|^3.
\end{equation}
Its threshold is not numerical. That theorem adapts the uniform smoothing,
contamination and local comparison methods of He and Cheng
\cite{HeCheng2026} to nonidentically distributed arrays. Its finite input
is the common-diameter Bernoulli theorem of \cite{BellFinite2026}.
The present paper concerns the remaining extremal problem at positive
Lyapunov ratio.

\begin{definition}\label{def:representation}
An admissible representation consists of independent centered variables
$X_1,X_2,\ldots$, an independent centered infinitely divisible variable
$Y$, and a finite number
\begin{align}
 S&=\sum_{i\ge1}X_i+Y,\notag\\
 \log\E e^{itY}
   &=-\frac{at^2}{2}
     +\int_{\R\setminus\{0\}}(e^{ity}-1-ity)\,\nu(dy),
                                                        \label{eq:LK}\\
 \sum_i\Var(X_i)+a+\int y^2\nu(dy)&=1,\qquad
 \beta=\sum_i\E|X_i|^3+\int|y|^3\nu(dy)<\infty.         \label{eq:profile-moments}
\end{align}
Here $a\ge0$, the series converges in $L^2$, and $\nu$ is a L\'evy
measure with finite second moment. The number $\beta$ is the additive
third-moment quantity of the representation. A representation is
maximizing if $\beta>0$ and $\Delta(S)=\ci\beta$.
\end{definition}

The Gaussian term contributes zero to $\beta$. Indeed a Gaussian variable
of variance $a$ can be approximated by sums whose variance is $a$ and
whose third-moment sums tend to zero. Thus $\beta$ is not the third
absolute moment of $S$, and it is not determined by the law of $S$ alone.
The approximation statement for the whole representation is proved in
Section~\ref{sec:compactness}.

\begin{theorem}\label{thm:structure}
If $\ci>\ce$, there is a maximizing admissible representation whose
variables $X_i$ have at most three support points. Every maximizing
representation with this support property satisfies
\begin{equation}\label{eq:main-structure}
 a=0,\qquad \nu=0,\qquad
 \max\{|u|:u\in\supp X_i\}\longrightarrow0.
\end{equation}
If infinitely many $X_i$ are nonconstant, then the law of $S$ is
nonatomic and
\[
 \sum_i\bigl(1-\max_u\Pp(X_i=u)\bigr)=\infty.
\]
\end{theorem}

Reflection, if needed, gives a threshold $z$ at which
\begin{equation}\label{eq:max-data}
 F(z)-\Phi(z)=D=C\beta>0,
 \qquad C=\ci,\qquad A=\frac{3D+z\phi(z)}2,
 \qquad F(x)=\Pp(S\le x).
\end{equation}
All CDFs in this paper are closed CDFs. A left limit is written explicitly.
The function $F-\Phi$ is upper semicontinuous and tends to zero at both
ends of the line, so a positive supremum on this branch is attained.
The continuous normal CDF allows reflection to interchange the positive
branch with the branch $\Phi(x)-\Pp(S<x)$.

There is also a restriction on support configurations. For a finite set
$I$ of indices, put $T_I=S-\sum_{i\in I}X_i$ and
\begin{equation}\label{eq:completion}
 \mathcal A_I=
 \left\{(x_i)_{i\in I}\in\prod_{i\in I}\supp X_i:
       z-\sum_{i\in I}x_i\in\supp T_I\right\}.
\end{equation}
The support on the right is the closed support of the actual independent
remainder. This definition does not condition on the event $S=z$.

\begin{theorem}\label{thm:rank}
Fix a maximizing representation as in Theorem~\ref{thm:structure} and a
threshold satisfying \eqref{eq:max-data}. There is $k_0$ such that, for
every finite $I\subset\{k_0,k_0+1,\ldots\}$, the following holds:
if functions $f_i:\supp X_i\to\R$ satisfy
\[
             \sum_{i\in I}f_i(x_i)=c
             \quad\text{for all }(x_i)\in\mathcal A_I,
\]
then every $f_i$ is constant. Equivalently, choose one state $u_i^0\in\supp X_i$ for each $i$.
The indicator vectors
$\bigl(\ind_{\{x_i=u\}}\bigr)_{i\in I,\,u\in\supp X_i\setminus\{u_i^0\}}$
associated with $\mathcal A_I$ have full affine dimension
$\sum_{i\in I}(|\supp X_i|-1)$.
\end{theorem}

\begin{corollary}\label{cor:separated}
Such a representation cannot consist of a finite-support independent
prefix followed by infinitely many nonconstant finite blocks of
finite-support summands for which
\[
 H_k=\sum_{i\text{ in block }k}\diam(\supp X_i),\qquad
 \delta_k=\min_{u\ne v\in\supp(\sum_{i\text{ in block }k}X_i)}|u-v|
\]
satisfy
\begin{equation}\label{eq:separation}
              \sum_k H_k<\infty,
              \qquad \delta_k\ge\sum_{l>k}H_l.
\end{equation}
There are no restrictions on the block sizes, their positive atom masses,
or the overlap of the prefix translates.
\end{corollary}

Theorem~\ref{thm:structure} leaves both finite representations and
nonatomic infinite convolutions with shrinking supports. Such infinite
convolutions need not have a density. Theorem~\ref{thm:rank} is automatic
when every relevant remainder has support $\R$; it therefore does not
supply the missing regularity assertion.

The compactification uses the classical accompanying Poisson-law method
and L\'evy--Khintchine representation \cite{Levy1937};
Section~\ref{sec:compactness} proves the needed moment bounds.
Coordinate contamination follows He and Cheng
\cite[Section~4]{HeCheng2026}; the independent-coordinate derivatives,
L\'evy-intensity variations and contact propagation are derived below.
The Gaussian estimates use the zero-bias transformation of Goldstein and
Reinert \cite{GoldsteinReinert1997}, the transport estimate independently
obtained by Goldstein and Tyurin \cite{Goldstein2010,Tyurin2009}, and the
compound-Poisson Berry--Esseen bound of Korolev and Shevtsova
\cite{KorolevShevtsova2012}. The atom argument uses Sperner's theorem
\cite{Sperner1928}. These inputs are stated where they enter the proofs.

\section{Compactification and attainment}\label{sec:compactness}

\begin{lemma}\label{lem:threepoint}
For every $c>0$, a finite independent variance-one array satisfying
$\Delta(\sum_iX_i)>c\sum_i\E|X_i|^3$ can be replaced by one with the
same number and variances of summands, each having at most three support
points, for which the same strict inequality holds.
\end{lemma}
\begin{proof}
Choose the positive closed-CDF branch by reflection. At its fixed threshold,
with all other summands held fixed, maximize
\[
       \int\{G_i(z-u)-c|u|^3\}\,P(du)
       \quad\text{subject to }\int u\,dP=0,\quad\int u^2\,dP=v_i.
\]
For $v_i=0$ the law is fixed. Otherwise the objective is upper
semicontinuous and bounded above by $1-c|u|^3$. A maximizing sequence
has bounded third moments, so a weak limit preserves its first two
moments and attains the supremum. The maximizing set is compact and is
a face of the convex moment set. It has an extreme point. Such a point
has at most three atoms: four disjoint positive-mass sets would give a
linear dependence among their mass, first-moment and second-moment
vectors, and hence two distinct feasible bounded-density perturbations
with the original law as their average. This is the elementary
three-constraint case of the moment-set principle in
\citet{Winkler1988}; the coercive argument is also given in
\citet{BellFinite2026}. Replacing coordinates successively preserves the
strict positive excess at coefficient $c$.
\end{proof}

\begin{lemma}\label{lem:compactness}
Let finite independent centered arrays have total variance one and
third-moment sums $\beta_n\le M$. After passage to a subsequence, their
sums converge weakly to an admissible representation with
\[
                      \beta\le\liminf_n\beta_n.
\]
If the original summands have at most three atoms, so do the limiting
independent variables $X_i$.
\end{lemma}
\begin{proof}
First take a subsequence along which $\beta_n$ tends to its liminf.
Order each row by decreasing variance, padding it with zero variables.
Then $v_{n,k}\le1/k$. The laws are tight, and
\[
 \E[X_{n,k}^2;|X_{n,k}|>L]\le M/L
\]
makes their squares uniformly integrable. A diagonal subsequence gives
$X_{n,k}\Rightarrow X_k$, convergence of means and variances, and
convergence of the numbers $\E|X_{n,k}|^3$ to numbers $b_k$.
Writing $v_k=\E X_k^2$ and $r_k=\E|X_k|^3$, we have
\[
 \sum_kv_k\le1,\qquad r_k\le b_k,
 \qquad\sum_kb_k\le\lim_n\beta_n.
\]
The independent series $\sum_kX_k$ exists in $L^2$.

Choose $K_n\to\infty$ sufficiently slowly that the first $K_n$ laws
approximate their limits on each fixed prefix and
$\sum_{k\le K_n}|v_{n,k}-v_k|\to0$. The corresponding prefix sums
converge weakly to $\sum_kX_k$: after fixing a finite prefix, the
variance of what remains in the growing prefix is at most
$\sum_{k>K}v_k+o(1)$. Define the quadratic measures of the remaining row by
\[
 \Lambda_n(E)=\sum_{k>K_n}\E[X_{n,k}^2\ind_{\{X_{n,k}\in E\}}].
\]
Their masses tend to $v_0=1-\sum_kv_k$ and their tails are bounded by
$M/L$ outside $[-L,L]$. Take $\Lambda_n\Rightarrow\Lambda$; its mass
is $v_0$.

For a centered variable of variance $v$, its characteristic function
$f$ satisfies $|f(t)-1|\le vt^2/2$. If $\Re w\le0$, then
$|e^w-1-w|\le|w|^2/2$. Telescoping products therefore gives
\begin{equation}\label{eq:poissonization}
 \left|\prod_jf_j(t)-\exp\Bigl\{\sum_j(f_j(t)-1)\Bigr\}\right|
                  \le\frac{t^4}{8}\sum_jv_j^2.
\end{equation}
For the remaining row the right side tends to zero, since its largest
variance is at most $1/(K_n+1)$. Its accompanying Poisson exponent is
\[
 \int q_t(y)\,\Lambda_n(dy),\qquad
 q_t(y)=\frac{e^{ity}-1-ity}{y^2},\quad q_t(0)=-t^2/2.
\]
The function $q_t$ is bounded and continuous. Hence these exponents
converge to the integral against $\Lambda$. Set
$a=\Lambda(\{0\})$ and $\nu(dy)=y^{-2}\Lambda(dy)$ off zero. This
gives \eqref{eq:LK}, with variance $v_0$. Independence of the original
prefix and tail proves the claimed representation and preserves total
variance.

For each fixed $K$ and all sufficiently large $n$,
\[
 \int|y|\,\Lambda_n(dy)
       \le\beta_n-\sum_{k\le K}\E|X_{n,k}|^3.
\]
Lower semicontinuity and then $K\to\infty$ give
$\int|y|^3\nu(dy)\le\lim\beta_n-\sum_kb_k$.
Together with $r_k\le b_k$ this proves the third-moment inequality.
No uniform integrability of cubic moments was assumed. Finally, a weak
limit of three-atom laws cannot charge four disjoint open sets.
\end{proof}

\begin{lemma}\label{lem:realization}
Every admissible representation is a weak limit of finite centered
independent arrays with variance exactly one and third-moment sums tending
to its number $\beta$. Consequently
\begin{equation}\label{eq:extended-bound}
                         \Delta(S)\le\ci\beta.
\end{equation}
\end{lemma}
\begin{proof}
Keep finitely many $X_i$. Truncate $\nu$ away from zero and infinity,
and approximate it by a finite atomic measure with convergent second
and third absolute integrals. A small contraction can make its variance
no larger than the available variance. Replace each atom
$\lambda\delta_y$ by $N$ independent variables $y(B-\lambda/N)$,
where $B\sim\Bern(\lambda/N)$. Their total variance and third-moment
sum are
\[
 \lambda y^2(1-\lambda/N),\qquad
 \lambda|y|^3(1-\lambda/N)
       \{(\lambda/N)^2+(1-\lambda/N)^2\}.
\]
They converge to the corresponding compensated Poisson law and its
second and third jump integrals. Fill the remaining variance by $N$
independent symmetric two-point variables. If that variance is $b$, their
third-moment sum is $b^{3/2}/\sqrt N$. Sending the truncation and
approximation indices to their limits gives the required Gaussian
variance and discards a tail of $\sum_iX_i$ whose variance tends to zero.
Every finite array has variance one.

If $P_n\Rightarrow P$ and $H$ is a continuous CDF, then
$\|F_{P_n}-H\|_\infty\to\|F_P-H\|_\infty$.
To verify this assertion, sandwich CDF values between a finite grid of
continuity points of $F_P$ on which the oscillation of $H$ is arbitrarily
small, and control both tails. Values immediately before and after jumps
are recovered by continuity points and the continuity of $H$.
Applying this to $H=\Phi$ proves \eqref{eq:extended-bound}.
\end{proof}

\begin{proposition}\label{prop:attainment}
If $\ci>\ce$, a maximizing representation with three-point independent
variables exists.
\end{proposition}
\begin{proof}
Take coefficients $c_n\uparrow\ci$ with $c_n>\ce$ and finite arrays
whose ratios exceed $c_n$. Lemma~\ref{lem:threepoint} supplies three-point
arrays with the same property. Their third-moment sums satisfy
$\ell_0<\beta_n<1/c_n$ by \eqref{eq:small-input} and $\Delta\le1$.
Take $\beta_n\to r>0$ and apply Lemma~\ref{lem:compactness}. Weak
continuity of distance to the normal gives $\Delta(S)=\ci r>0$.
The limiting representation satisfies
\[
          \ci r=\Delta(S)\le\ci\beta\le\ci r.
\]
Thus $\beta=r$ and it is maximizing. This equality rules out loss of the
additive third-moment quantity in the compactification.
\end{proof}

\section{Variations and contact identities}\label{sec:variation}
Fix a maximizing representation with $C=\ci>\ce$ and
\eqref{eq:max-data}. For an admissible perturbation of raw independent
laws, let $m$ be the total mean, $B$ the total central variance, $R_*$
the sum of the recentered third absolute moments plus the L\'evy third
integral, and $F_*$ the raw sum CDF. At a raw threshold $w$, put
\begin{equation}\label{eq:objective}
       J=B^{3/2}\left\{F_*(w)-\Phi\left(\frac{w-m}{\sqrt B}\right)\right\}
                    -C R_*.
\end{equation}
By Lemma~\ref{lem:realization}, $J\le0$ whenever the perturbation is
admissible, and $J=0$ at the original representation and $w=z$.
Finite probability variations make $F_*(z)$ a multiaffine polynomial,
including at atoms. Since $|x|^3$ is $C^2$, their central third moments
are twice differentiable. A two-sided probability parameter consequently
has zero first derivative, and the Hessian is negative semidefinite on
its admissible tangent cone when all first derivatives vanish.

\begin{lemma}\label{lem:contacts}
Let $G_i$ be the CDF of $S-X_i$, and put
$v_i=\E X_i^2$, $b_i=\E|X_i|^3$, $M_i=\E[X_i|X_i|]$. Then
\begin{align}
 G_i(z-u)-F(z)+\phi(z)u+A(u^2-v_i)
   -C\{|u|^3-b_i-3M_i u\}&\le0,                    \label{eq:coord-contact}\\
 Q(y):=F(z-y)-F(z)+\phi(z)y+Ay^2-C|y|^3&\le0.       \label{eq:Q}
\end{align}
Equality holds in \eqref{eq:coord-contact} at every support point of
$X_i$, and $Q(y)=0$ at every nonzero point of $\supp\nu$.
All supports $\supp X_i$ lie in one bounded interval, and all nonzero
zeros of $Q$ lie in one bounded interval.
\end{lemma}
\begin{proof}
Replace the raw law of $X_i$ by $(1-t)\Law(X_i)+t\delta_u$.
The derivatives of mean, variance and third absolute central moment are
$u$, $u^2-v_i$ and $|u|^3-b_i-3M_i u$. Differentiating
\eqref{eq:objective} gives \eqref{eq:coord-contact}. Its integral under
$\Law(X_i)$ is zero. Upper semicontinuity then extends equality from
almost every point to the entire support. This is the coordinatewise
contamination argument of \citet[Section~4]{HeCheng2026}.

Append a compensated Poisson variable $y(N_t-t)$ and use the centered
threshold $z-ty$. Its raw CDF is exactly $\E F(z-yN_t)$; variance and
third-moment quantity increase by $ty^2$ and $t|y|^3$. Its right
first derivative is $Q(y)$, proving \eqref{eq:Q}. For a finite signed
measure $\eta=f\nu$, with bounded $f$ supported away from zero, both
signs of small changes in $\nu$ are admissible. Their derivatives are
$\int Q\,d\eta$. Thus $Q=0$ almost everywhere for $\nu$, and upper
semicontinuity gives equality on its nonzero support.

In \eqref{eq:coord-contact}, use $v_i\le1$, $|M_i|\le1$ and
$Cb_i\le C\beta=D$. Its right side at a support point is bounded
above by a fixed cubic polynomial in $|u|$ with leading coefficient
$-C$. The same reasoning applies to \eqref{eq:Q}. Both zero sets are
therefore bounded uniformly in $i$.
\end{proof}

For later calculations, consider a mass transfer between two states
$u,v$ of one centered variable. Write $d=u-v$, $b=u^2-v^2$, and let
$g=G_i(z-u)-G_i(z-v)$. The third-moment derivatives are
\begin{align}
 r'&=|u|^3-|v|^3-3d\E[X_i|X_i|],\notag\\
 r''&=-6d(u|u|-v|v|)+6d^2\E|X_i|.                  \label{eq:transfer-moments}
\end{align}
Stationarity says $g+\phi d+Ab-Cr'=0$. Hereafter $\phi$ without an
argument in an algebraic formula means $\phi(z)$. Put
\[
 B_0=\frac{3D+(3z+z^3)\phi}{4},\qquad
 E_0=(1+z^2/2)\phi.
\]
For two transfers in distinct coordinates the Hessian entries of
\eqref{eq:objective} are
\begin{align}
 H_{ii}&=3b_i g_i+B_0b_i^2+2E_0b_i d_i-3Dd_i^2-Cr_i'',\label{eq:prob-hessian}\\
 H_{ij}&=g_{ij}+\tfrac32(b_i g_j+b_j g_i)+B_0b_i b_j
      +E_0(b_i d_j+b_j d_i)+z\phi d_i d_j,\quad i\ne j.\notag
\end{align}
Here $g_{ij}$ is the exact raw double CDF difference. These identities
follow by differentiating $B^{3/2}\Phi((z-m)/\sqrt B)$ at $(B,m)=(1,0)$;
its second derivatives are
\[
 N_{mm}=-z\phi,\quad N_{mB}=-(1+z^2/2)\phi,\quad
 N_{BB}=\tfrac34\Phi(z)-\tfrac14(3z+z^3)\phi.
\]
The second variance derivative in one transfer is $-2d_i^2$, whereas
mixed moment derivatives in different coordinates vanish. This accounts
for the diagonal term $-3Dd_i^2$ in \eqref{eq:prob-hessian}.

\begin{lemma}\label{lem:atoms-modal}
If the full sum has an atom of mass $m>0$ and each $X_i$ has at most
three support points, then, with $q_i=\max_u\Pp(X_i=u)$,
\[
                         \sum_i(1-q_i)<\infty.
\]
\end{lemma}
\begin{proof}
Represent each $X_i$ as an activated fair two-point law. If $q_i\ge1/2$,
activate with probability $\rho_i=2(1-q_i)$ and use the modal
state with the remaining probability. On activation choose fairly between the modal state and an
independent conditional nonmodal state. If three masses $p_1,p_2,p_3$
are all below $1/2$, mix the three fair pairs, assigning weight $1-2p_k$
to the pair that excludes state $k$. Thus in all cases
$\rho_i=\min\{1,2(1-q_i)\}\ge\tfrac32(1-q_i)$.

Condition on activations and chosen endpoints, orienting the fair bits
with positive increments. Every fixed-sum fiber is an antichain.
Sperner's theorem \cite{Sperner1928} bounds its conditional probability
by $2^{-N}\binom{N}{\lfloor N/2\rfloor}\le\sqrt{2/(N+1)}$ when
$N$ bits are active. For a finite prefix with $\mu=\sum_i\rho_i>0$,
Cauchy--Schwarz and the generating function of $N$ yield
\[
 \sup_x\Pp(\textstyle\sum_iX_i=x)^2
 \le2\E\frac1{N+1}
 =2\int_0^1\prod_i(1-\rho_i+\rho_i t)\,dt
 \le\frac{2(1-e^{-\mu})}{\mu}.
\]
Convolution cannot increase the largest atom. The left side is at least
$m^2$ for every prefix, so $\mu$ stays bounded, proving the assertion.
This is a direct finite-alphabet concentration argument within the
classical theory of infinite convolutions \cite{JessenWintner1935}.
\end{proof}

\section{Eliminating jumps and nonshrinking supports}\label{sec:propagation}

\subsection{Pair contacts for the L\'evy measure}

\begin{lemma}\label{lem:levy-pairs}
For nonzero same-sign points $x,y\in\supp\nu$, one has $Q(x+y)=0$.
The measure $\nu$ is supported on at most two positive and two negative
points.
\end{lemma}
\begin{proof}
Choose a probability law $H$ dominated by a finite multiple of $\nu$
and supported in a compact interval on one side of zero. Write
$m=\E H$, $v=\E H^2$, $r=\E|H|^3$, and $k=m^2/v$. Remove a
compensated Poisson component of intensity $p-kp^2$ with jump law $H$,
and insert the single centered variable $B_pH-pm$. This is admissible
for all sufficiently small $p>0$ and preserves variance. Its third
absolute moment is
\[
      pr-3p^2|m|v+4p^3|m|^3-2p^4|m|^3.
\]
At the centered threshold $z-kp^2m$, the raw CDF is
\[
       [(I+pU_H)e^{-(p-kp^2)U_H}F](z),\qquad
       U_H f(w)=\E f(w-H)-f(w).
\]
The bounded-operator expansion is valid in uniform norm. Its second-order
term is $(kU_H-\tfrac12U_H^2)F$. Combining it with the normal and
third-moment terms, the coefficient of $p^2$ in the excess is
\[
                 (k+1)\E Q(H)-\tfrac12\E Q(H+H'),
\]
where $H'$ is independent. The cubic cancellation uses
$\E|H+H'|^3=2r+6|m|v$. The first expectation vanishes by
Lemma~\ref{lem:contacts}; the second is nonpositive. Maximality forces
it to be zero. Upper semicontinuity extends this identity to all
same-sign support pairs, by choosing $H$ on neighborhoods of any two
specified support points. The case $x=y$ is included.

For the support count, perturb the L\'evy measure by a finite signed
measure $\eta=f\nu$ with bounded $f$ supported away from zero. Define
$d_\eta=\int y\,d\eta$, $v_\eta=\int y^2\,d\eta$, and
$r_\eta=\int|y|^3\,d\eta$. The centered threshold $z-d_\eta$ gives
\begin{equation}\label{eq:intensity-objective}
 (1+v_\eta)^{3/2}
 \left\{[e^{U_\eta}F](z)
       -\Phi\left(\frac{z-d_\eta}{\sqrt{1+v_\eta}}\right)\right\}
        -C(\beta+r_\eta)\le0,
\end{equation}
where $U_\eta f(w)=\int[f(w-y)-f(w)]\,\eta(dy)$.
Both signs are allowed near zero. Negative perturbations are interpreted
by removing the corresponding finite Poisson component; the bounded
operator identity remains exact. Thus its Hessian is negative
semidefinite on these signed directions.

Write $U_xf(w)=f(w-x)-f(w)$. Let $L=(15D+\phi'''(z))/4$, and, for $\sigma\in\{-1,1\}$, define
\begin{equation}\label{eq:kernel}
 \mathcal K_\sigma(u,v)=uv\left\{-3D+
       \left(3C+\frac{\sigma\phi''(z)}2\right)(u+v)
       -Luv+\frac{3C}{2}uv(u+v)\right\},\quad u,v>0.
\end{equation}
At same-sign contact points $x=\sigma u$, $y=\sigma v$ whose sum is
also a contact, the intensity Hessian equals $\mathcal K_\sigma(u,v)$.
For verification before substituting contacts, its kernel is
\begin{align*}
 U_xU_yF(z)&+\tfrac32[x^2U_yF(z)+y^2U_xF(z)]
      +B_0x^2y^2\notag\\[-1ex]
 &+E_0(x^2y+xy^2)+z\phi xy.
\end{align*}
Substitute $U_xF(z)=-\phi x-Ax^2+C|x|^3$ and
$U_xU_yF(z)=-2Axy+3Cxy(|x|+|y|)$ to obtain
\eqref{eq:kernel}. Its moment coefficient matrix in the basis
$(u,u^2,u^3)$ is
\begin{equation}\label{eq:moment-matrix}
 M_\sigma=
 \begin{pmatrix}
 -3D&3C+\sigma\phi''(z)/2&0\\
 3C+\sigma\phi''(z)/2&-L&3C/2\\
 0&3C/2&0
 \end{pmatrix}.
\end{equation}
It is not negative semidefinite: its value at $(0,t,1)$ is
$-Lt^2+3Ct>0$ for small $t>0$.

If three distinct support magnitudes $u_1,u_2,u_3$ existed on one
side, restrict $\nu$ to three disjoint small compact neighborhoods and
normalize each restriction to a probability measure. Arbitrary signed
combinations are locally admissible intensity directions. Their Hessian
is congruent to $M_\sigma$ through the columns of first, second and
third moments of these restrictions. As the neighborhoods shrink, these
columns tend to $(u_j,u_j^2,u_j^3)$, an invertible Vandermonde matrix.
For sufficiently small neighborhoods the Hessian cannot be negative
semidefinite. This contradiction proves the support count. Finite second
moment then makes $\nu$ a finite measure.
\end{proof}

\subsection{Contacts from limiting support values}

\begin{lemma}\label{lem:limit-double}
Suppose $i_n$ are distinct and $u_n\in\supp X_{i_n}$ satisfies
$u_n\to x\ne0$. Then $Q(x)=Q(2x)=0$. The probability transfer can be
taken from a state $v_n\to0$ to $u_n$, with gap $d_n=u_n-v_n\to x$.
If the full CDF has a jump and $d_n>x$ along a subsequence, it has no
jump at $z-x$ or $z-2x$.
\end{lemma}
\begin{proof}
The variances of distinct coordinates tend to zero. If $F$ is continuous,
choose a support state $v_n$ with $|v_n|\le\sqrt{v_{i_n}}$. Deleting
one or two selected coordinates changes the sum by a variable tending
to zero in $L^2$. The deleted CDFs therefore converge uniformly to $F$.
The derivatives in \eqref{eq:transfer-moments} satisfy
\[
 d_n\to x,\quad b_n=u_n^2-v_n^2\to x^2,
 \quad r_n'\to|x|^3,\quad r_n''\to-6|x|^3.
\]
First stationarity gives $Q(x)=0$. In
\eqref{eq:prob-hessian}, the limit of the off-diagonal entry minus
that of the diagonal entry is
\[
 F(z-2x)-2F(z-x)+F(z)+2Ax^2-6C|x|^3=Q(2x).
\]
The two-coordinate contrast $(1,-1)$ gives $-2Q(2x)\le0$.
Together with $Q\le0$, this proves the assertion.

If the full law has an atom, Lemma~\ref{lem:atoms-modal} gives modal
probabilities $q_n\to1$. Take $v_n$ to be a modal state; then
$q_nv_n^2\le v_{i_n}$, so $v_n\to0$. Put
$H_n(w)=\Pp(S-X_{i_n}+v_n\le w)$. Since the original law contains
this branch with probability $q_n$,
\[
                    \|H_n-F\|_\infty\le1-q_n.
\]
The analogous double-deletion bound is $(1-q_n)+(1-q_m)$.
Pass to a subsequence on which the gaps equal $x$, lie above $x$,
or lie below $x$. Each shifted CDF then has a definite side limit.
Let $Q_1^*,Q_2^*$ denote the corresponding versions of $Q(x),Q(2x)$,
using these side limits. Stationarity gives $Q_1^*=0$; the same Hessian
contrast gives $Q_2^*\ge0$, while $Q\le0$ gives $Q_2^*\le0$.
Closed CDF values dominate their left limits, so $Q(x),Q(2x)\ge0$.
Both are zero. When the arguments approach from the left, the two side
limits must equal the closed values; this proves the last assertion.
\end{proof}

\subsection{Comparisons with fixed added summands}
Fix a sign $\sigma$ and a nonzero limit $x=\sigma u$ as in
Lemma~\ref{lem:limit-double}. Its probability transfers are two-sided
at each finite index. Appending a fresh raw variable
$(1-t)\delta_0+t\delta_y$, $t\ge0$, has first derivative $Q(y)$
in \eqref{eq:objective}. The gap $y$ below is always fixed, rather than
set equal to a varying support value. Thus when $Q(y)=0$, its first
derivative is exactly zero for every selected index.

\begin{lemma}\label{lem:closed-copy}
For $y=\sigma v$ with $Q(y)=0$, the limiting mixed Hessian between the
old transfer tending to $x$ and a fresh variable at $y$ equals
\[
                     \mathcal K_\sigma(u,v)+Q(x+y).
\]
The old diagonal and its mixed entry with a fresh variable at $x$ both
equal $\mathcal K_\sigma(u,u)$ in the limit.
The same conclusions hold if the old direction instead varies the
positive intensity of an existing L\'evy atom at $x$.
\end{lemma}
\begin{proof}
The limiting derivatives of an old probability transfer are
\[
 m'=x,\quad B'=x^2,\quad B''=-2x^2,\quad
 R_*'=|x|^3,\quad R_*''=-6|x|^3.
\]
They coincide with the derivatives for the fresh variable at $x$.
Stationarity and \eqref{eq:prob-hessian} give the common limiting
diagonal $\mathcal K_\sigma(u,u)$. Their mixed CDF double shift is
closed: in the continuous case this follows from uniform deletion
convergence; in the atomic case its argument approaches $z-2x$ from
the right, or Lemma~\ref{lem:limit-double} removes the left-limit jump.
Since $Q(x)=Q(2x)=0$, its mixed entry is the same diagonal.

Append also a fresh variable at the fixed contact $y$. Pass to a
subsequence of the bounded Hessian entries. Write $B_{xy}$ for the
old--fresh-$y$ entry and $N_{xy}$ for the fresh-$x$--fresh-$y$ entry.
The latter uses the exact closed value $F(z-x-y)$ and, by the same
substitution as in \eqref{eq:kernel}, equals
$\mathcal K_\sigma(u,v)+Q(x+y)$. The old single-shift limit is fixed
by stationarity; the fresh single shift is closed. Upper semicontinuity
of the double shift, together with the uniform deletion or modal-branch
comparison in the preceding lemma, therefore gives $N_{xy}-B_{xy}\ge0$.

The vector $w=(-1,1,0)$ has zero limiting quadratic value and is feasible:
it decreases an old probability and increases its fresh copy. For every
$t\ge0$, $w+t(0,0,1)$ is also a feasible tangent direction. The Hessian
inequality gives
\[
                   2t(N_{xy}-B_{xy})+t^2H_{yy}\le0.
\]
Letting $t\downarrow0$ proves the reverse inequality. This identifies
the closed mixed entry. All first derivatives were exactly zero before
taking the limit, so no approximation of a first-order contact is used
to infer a second-order inequality.

For an existing L\'evy atom, vary its intensity in both directions. Its
CDF and moment derivatives are already at fixed closed shifts. Relative
to the fresh Bernoulli diagonal its second derivative differs by
\[
             U_x^2F(z)+2Ax^2-6C|x|^3=Q(2x)-2Q(x)=0,
\]
using Lemma~\ref{lem:levy-pairs}. This proves the same formulas.
\end{proof}

\begin{proposition}\label{prop:shrink}
Every maximizing representation with three-point independent variables
has $\nu=0$ and $\max\{|u|:u\in\supp X_i\}\to0$.
\end{proposition}
\begin{proof}
Let $x=\sigma u$ be either a nonzero limiting support value or a point
of $\supp\nu$. In the latter case it is an atom of positive intensity
by Lemma~\ref{lem:levy-pairs}. Let $y=\sigma v$ be any same-sign global
contact. For fixed $\eps>0$, append the one raw three-point law
\[
        (1-(1+\eps)t)\delta_0+t\delta_x+\eps t\delta_y,
                      \qquad t\ge0.
\]
Its first derivative is $Q(x)+\eps Q(y)=0$. If
$m_1=x+\eps y$, $v_1=x^2+\eps y^2$, and
$s_1=x|x|+\eps y|y|$, its moment derivatives are
\[
 m'=m_1,\quad B'=v_1,\quad B''=-2m_1^2,
 \quad R_*'=|x|^3+\eps|y|^3,\quad R_*''=-6m_1s_1.
\]
The raw CDF is affine in $t$. Its Hessian diagonal is consequently
\[
 \mathcal K_\sigma(u,u)+2\eps\mathcal K_\sigma(u,v)
                              +\eps^2\mathcal K_\sigma(v,v).
\]
By Lemma~\ref{lem:closed-copy}, its mixed entry with the old direction
is $\mathcal K_\sigma(u,u)+
\eps[\mathcal K_\sigma(u,v)+Q(x+y)]$.
Evaluate the two-parameter limiting Hessian on $(-1,1)$, decreasing
the old direction and increasing the new activation. It gives
\[
                  -2\eps Q(x+y)+\eps^2\mathcal K_\sigma(v,v)\le0.
\]
First take the old-coordinate limit with $\eps$ fixed, and only then
let $\eps\downarrow0$. Thus $Q(x+y)\ge0$, and hence $Q(x+y)=0$.
The admissible neighborhoods of the old probabilities need not have a
uniform radius.

Starting with $Q(x)=0$, this conclusion propagates to $Q(nx)=0$ for
every positive integer $n$, always using the same old direction and a
fixed contact at the next step. This contradicts the bounded contact
set in Lemma~\ref{lem:contacts}. No such $x$ exists. The finite L\'evy
measure must therefore be zero. The uniformly bounded independent
supports cannot have a nonzero accumulation along distinct indices, so
their radii tend to zero.
\end{proof}

\section{Excluding a Gaussian component}\label{sec:gaussian}
By Proposition~\ref{prop:shrink}, it is enough to consider
\[
 S=W+G_a,\qquad W=\sum_iX_i,\qquad
 \Var(W)=1-a,\quad\beta=\sum_i\E|X_i|^3,
\]
where $G_a\sim N(0,a)$ is independent. Suppose $a>0$. The law has a
smooth density $f$. Threshold and Gaussian-variance stationarity give
\begin{equation}\label{eq:gauss-stationarity}
 f(z)=\phi(z),\qquad f'(z)=-3D-z\phi(z),\qquad |f''(z)|\le6C.
\end{equation}
The second identity follows from the heat equation and the derivatives of
\eqref{eq:objective}. Both signs of a small variance change are allowed.
The last inequality follows by expanding $Q$ at zero on its two sides.

\subsection{A smoothing bound from zero bias}
For a centered variable $T$ of variance $v>0$, its zero-biased law $T^*$
is characterized by $\E[Tg(T)]=v\E g'(T^*)$.
Goldstein and Reinert \cite{GoldsteinReinert1997} introduced this
transformation and the independent-sum replacement construction.
Goldstein \cite[Lemma~1.1]{Goldstein2010} and Tyurin
\cite[Theorem~3 and Corollary~1]{Tyurin2009} give
\begin{equation}\label{eq:zero-bias}
                  W_1(W,W^*)\le\frac{\beta}{2(1-a)},
\end{equation}
where $W_1$ is Wasserstein distance with cost $|x-y|$. For countably many
summands the statements below follow by finite approximation retaining
the fixed Gaussian variance.

\begin{lemma}\label{lem:gauss-large}
For every such independent sum,
\[
 \Delta(W+G_a)\le K(a)\beta,\qquad
 K(a)=\frac{\phz}{2}\int_0^1
            \frac{s^2}{[1-(1-a)s^2]^{3/2}}\,ds.
\]
The function $K$ is decreasing, and $K(1/10)<\ce$. Hence a maximizing
representation has $a<1/10$.
\end{lemma}
\begin{proof}
For $h(u)=\Phi((z-u)/\sqrt a)$ and $v=1-a$, the solution of
$vg'(u)-ug(u)=h(u)-\E h(\sqrt v Z)$ has the
Ornstein--Uhlenbeck representation
\[
              g(u)=-\int_0^1\E h'(su+\sqrt{v(1-s^2)}Z)\,ds.
\]
Differentiating twice and convolving Gaussian derivatives gives
$\|g''\|_\infty\le\phz\int_0^1s^2(1-vs^2)^{-3/2}\,ds$.
The zero-bias identity and \eqref{eq:zero-bias} bound the expectation
difference by $\beta\|g''\|_\infty/2$.
The integral decreases with $a$, and direct integration gives
\[
 K(a)=\frac{\phz}{2}
 \left\{\frac1{(1-a)\sqrt a}
       -\frac{\arcsin\sqrt{1-a}}{(1-a)^{3/2}}\right\}.
\]
At $a=1/10$ this is
$(5\sqrt{10}/27)\phz(3-\arctan3)$.
To check the strict comparison without decimal rounding, use
$\pi>157/50$, $\sqrt{10}<19/6$ and
\[
 \arctan3=\pi/2-\arctan(1/3)
 >\frac{157}{100}-\left(\frac13-\frac1{81}+\frac1{1215}\right)
 >\frac{237}{190}
 >\frac{21}{10}-\frac{27}{10\sqrt{10}}.
\]
This is equivalent to $K(1/10)<\ce$.
\end{proof}

\subsection{A lattice comparison}

\begin{lemma}\label{lem:lattice-gauss}
At a maximizing representation with $a>0$,
\begin{equation}\label{eq:lattice-curvature}
             \phi(z)\le a\left(C-\frac{f''(z)^2}{36C}\right).
\end{equation}
\end{lemma}
\begin{proof}
We first approximate $W+G_a$ by a lattice convolution plus a Gaussian,
with error smaller than its lattice spacing. Order the $X_i$ by decreasing
variance and set
\[
 v^{(K)}=\sum_{i>K}v_i,\qquad H_K=\sum_{i>K}v_i^2.
\]
Then $KH_K\le v^{(K)}\to0$. If the series is infinite, choose
$h_K=\sqrt{H_K/K}$, so $H_K=o(h_K)$ and $Kh_K\to0$.
For finitely many nonzero summands take $h_K=K^{-2}$ beyond them.

Replace the tail after $K$ by its accompanying compensated Poisson law
with measure $\mu_K=\sum_{i>K}\Law(X_i)|_{\R\setminus\{0\}}$.
Variance and additive third moment are unchanged. By
\eqref{eq:poissonization}, after convolution with $G_a$ the uniform
CDF error is at most
\[
       \frac{H_K}{8\pi}\int_0^\infty t^3e^{-at^2/2}\,dt
                       =\frac{H_K}{4\pi a^2}=o(h_K).
\]
Extra powers of $t$ give the corresponding density-derivative estimates.

Round each of the first $K$ summands to its neighboring points on
$h_K\Z$, independently and with its conditional mean unchanged.
Writing the rounding error as $\delta$, we have
$\E(\delta\mid X)=0$, $|\delta|\le h_K$ and
$\E(\delta^2\mid X)\le h_K^2/4$. Subtract its total added variance
from the Gaussian component. Gaussian Taylor estimates bound the CDF
change by $O_a(Kh_K^2)=o(h_K)$ and likewise control fixed density
derivatives. The third-moment change is at most
\[
             3h_K^2\sum_{i\le K}\E|X_i|+Kh_K^3=o(h_K),
\]
since the sum of first absolute moments is at most $\sqrt K$.

Replace the jumps of $\mu_K$ smaller than $h_K$ in magnitude by their
Gaussian variance, and round all larger jumps without changing their
conditional means. Subtract their added variance from the Gaussian part.
For a large jump $y$, with rounding variance $e_h(y)$, the compensated
exponent error is bounded by
\[
 \left|\E e^{it(y+\delta)}-e^{ity}+\tfrac12t^2e_h(y)\right|
       \le |t|^3\left(\tfrac12h^2|y|+\tfrac16h^3\right).
\]
For a small jump it is at most $|t|^3|y|^3/6$. Integration against
$\mu_K$ gives exponent error $O(|t|^3h_Kv^{(K)})$ and third-moment
change $O(h_Kv^{(K)})=o(h_K)$. The Gaussian variance changes by
$o(1)$ and remains at least $a/2$. Factoring a common smaller Gaussian
factor, the bound $|e^u-e^v|\le|u-v|$ for $\Re u,\Re v\le0$
justifies the CDF and derivative comparisons by Fourier inversion.

We have obtained variance-one representations $T_j+G_{a_j}$, where
$T_j$ lies on a translate of $h_j\Z$, $a_j\to a$, and
\begin{equation}\label{eq:lattice-approx}
 \beta_j=\beta+o(h_j),\quad
 \|F_{T_j+G_{a_j}}-F\|_\infty=o(h_j),\quad
 f_j^{(k)}\to f^{(k)}\text{ uniformly},\quad k=0,1,2.
\end{equation}
The rounded jump measure has finite intensity, because the pre-rounded
jumps have magnitude at least $h_j$ and finite second integral. Its
centering is well-defined. All approximations are admissible representations.

Fix $p\in(0,1)$. Choose integers $n_j$ and $p_j\to p$ so that
$n_jh_j^2p_j(1-p_j)=a_j$; rounding $a_j/[h_j^2p(1-p)]$ to a nearby
integer and solving for $p_j$ does this, with the upper integer used
when $p=1/2$. Put $u_j=1-2p_j$ and
$B_j=h_j(\Bin(n_j,p_j)-n_jp_j)$. The lattice Edgeworth expansion at
closed endpoints is
\begin{equation}\label{eq:binomial-expansion}
 \Pp(B_j\le x)=\Phi_{a_j}(x)+\frac{h_j}{2}\phi_{a_j}(x)
              -\frac{a_jh_ju_j}{6}\phi_{a_j}''(x)+o(h_j),
\end{equation}
uniformly in those endpoints. Here $\Phi_b,\phi_b$ denote the centered
normal CDF and density of variance $b$. This is the binomial case of the
classical lattice expansion of Esseen; the closed-endpoint convention is
recorded, for example, in \citet[equation~(3)]{DinevMattner2013}.
For the uniformity with $p_j\to p$, Stirling's formula gives, with
$\sigma_j=\sqrt{n_jp_j(1-p_j)}$ and
$t_k=(k-n_jp_j)/\sigma_j$,
\[
 \Pp(\Bin(n_j,p_j)=k)=\frac{\phi(t_k)}{\sigma_j}
 \left[1+\frac{u_j(t_k^3-3t_k)}{6\sigma_j}
       +O\left(\frac{1+|t_k|^6}{n_j}\right)\right]
\]
uniformly for $|t_k|\le\sqrt{B\log n_j}$ and $p_j$ in a fixed
compact subinterval. Chernoff bounds control the omitted tails by
$o(n_j^{-1/2})$ when $B$ is sufficiently large. Summation of the remainder
is $O((\log n_j)^3/n_j)$, and right-endpoint Euler--Maclaurin summation
adds $\phi(t_k)/(2\sigma_j)$. Integration of the cubic Hermite term
is $(1-t_k^2)\phi(t_k)$, proving \eqref{eq:binomial-expansion}.

Choose $z_j$ on the full lattice of $T_j+B_j$ with $|z_j-z|\le h_j/2$.
Conditional on every value of $T_j$, the endpoint $z_j-T_j$ is a closed
binomial endpoint, so \eqref{eq:binomial-expansion} gives
\[
 F_{T_j+B_j}(z_j)=F_{T_j+G_{a_j}}(z_j)
         +\frac{h_j}{2}f_j(z_j)
         -\frac{a_jh_ju_j}{6}f_j''(z_j)+o(h_j).
\]
Threshold stationarity and \eqref{eq:lattice-approx} make its base
excess $D+o(h_j)$. The actual third-moment sum of the binomial carrier is
$a_jh_j(1+u_j^2)/2$. Maximality therefore implies, for $u=1-2p$,
\[
                    \phi(z)-\frac{au}{3}f''(z)-Ca(1+u^2)\le0.
\]
The inequality extends to $u\in[-1,1]$ by continuity. Optimize at
$u=-f''(z)/(6C)$, which lies in this interval by
\eqref{eq:gauss-stationarity}, to obtain
\eqref{eq:lattice-curvature}.
\end{proof}

Gaussian convolution gives the semiconvexity inequality
$f'(z)^2\le f(z)f''(z)+f(z)^2/a$ by Cauchy--Schwarz on its Gaussian
mixture formula. Set $c=\phi(z)/(aC)$. Lemma~\ref{lem:lattice-gauss}
gives $0<c\le1$ and $|f''(z)|\le6C\sqrt{1-c}$. Hence
\[
       (3D+z\phi(z))^2
       \le aC^2\{c^2+6c\sqrt{1-c}\}\le\frac{45}{16}aC^2.
\]
The last maximum is at $c=3/4$. Moreover,
$|z|\phi(z)\le\sqrt{2\phz/e}\sqrt{\phi(z)}$.
Since $C>2/5$, $\phz<2/5$ and $e>27/10$, it follows that
\[
 \beta\le\left(\frac{\sqrt5}{4}
        +\frac{\sqrt{2\phz/e}}{3\sqrt C}\right)\sqrt a
                   <\frac{17}{20}\sqrt a.
\]
For the final rational comparison one may use
$\sqrt5/4<14/25$ and $\sqrt{20/27}/3<29/100$.
Together with Lemma~\ref{lem:gauss-large}, this confines a possible
Gaussian component to
\begin{equation}\label{eq:gauss-window}
         0<\beta<\frac{269}{1000},\qquad
         \frac{400}{289}\beta^2<a<\frac1{10}.
\end{equation}

\subsection{Fourier exclusion of the remaining window}
Korolev and Shevtsova \cite[Section~3.3, Theorem~2]{KorolevShevtsova2012}
prove that a centered compound-Poisson variable of variance $v>0$ and
third absolute jump integral $r$ satisfies
\begin{equation}\label{eq:cp-bound}
                       d_K(Y,N(0,v))\le0.3051\,r/v^{3/2}.
\end{equation}
Their jump law is assumed to have positive variance. A vanishing symmetric
perturbation of a constant jump law removes this restriction. The same
bound holds for every finite-variance infinitely divisible law with third
jump integral $r$: truncate its L\'evy measure and replace the omitted
variance and Gaussian part by symmetric jumps of size $\eta$. These
add at most $v\eta$ to the third jump integral. Pass to the limit using
weak continuity of distance to the continuous normal CDF.

\begin{lemma}\label{lem:gauss-window-exclusion}
A variance-one independent sum satisfying \eqref{eq:gauss-window} has
$\Delta(S)<(2/5)\beta$.
\end{lemma}
\begin{proof}
The elementary inequality
\begin{equation}\label{eq:cosine}
                   \cos x\le1-x^2/2+|x|^3/10
\end{equation}
holds on the real line. For $0\le x\le12/5$ use the fourth-degree
cosine upper Taylor bound, and for $x\ge5$ its right side is at least
one. On $12/5\le x\le5$, the inequality
$\cos x\le-1+(x-\pi)^2/2$ reduces the claim to positivity of
$P(x)=2-x^2+\pi x-\pi^2/2+x^3/10$. With $u=x-4$,
\[
 P(x)=P(4)+(\pi-16/5)u+u^2(1/5+u/10).
\]
For $-8/5\le u\le0$ this is at least $P(4)$, and for $0\le u\le1$
it is at least $P(4)-(5/4)(\pi-16/5)^2$. Both are positive from
$157/50<\pi<22/7$. Evenness proves \eqref{eq:cosine}.

For independent centered copies $X,X'$, the pointwise inequality
\[
 |x-y|^3\le2(|x|^3+|y|^3)-2xy(|x|+|y|)
\]
gives $\E|X-X'|^3\le4\E|X|^3$. Thus a coordinate characteristic
function $h_i$ obeys
\begin{equation}\label{eq:cf-modulus}
 |h_i(t)|\le\exp\{-v_it^2/2+b_i|t|^3/5\},
          \qquad b_i=\E|X_i|^3.
\end{equation}
Write $h$ for the full characteristic function. On omitting coordinate
$i$, \eqref{eq:cf-modulus} and the exact Gaussian factor give
\[
 e^{-at^2/2}\prod_{j\ne i}|h_j(t)|
       \le e^{25/54-t^2/2+\beta t^3/5},\qquad t\ge0.
\]
Indeed $b_i\ge v_i^{3/2}$ and
$\max_{u\ge0}(u^2/2-u^3/5)=25/54$.
The single-coordinate form of \eqref{eq:zero-bias} gives
$|h_i(t)-h_i^*(t)|\le b_it/(2v_i)$. Selecting the replaced coordinate
with its variance probability in the zero-bias identity yields
\[
 |h'(t)+th(t)|\le\frac\beta2t^2
                    e^{25/54-t^2/2+\beta t^3/5}.
\]
Multiply by $e^{t^2/2}$ and integrate from zero. Then
\begin{equation}\label{eq:cf-difference}
 |h(t)-e^{-t^2/2}|
 \le\frac56 e^{25/54-t^2/2}(e^{\beta t^3/5}-1)
 \le\frac\beta6e^{25/54}t^3e^{-t^2/2+\beta t^3/5}.
\end{equation}
These calculations first apply to a finite sum. Finite approximations
retaining $a$ and converging in variance and third-moment sum justify
both inequalities for the countable sum.

Set $c=19/10$ and $T=c/\beta>7$. On $[0,6]$ use
$\beta\le269/1000$, and on $[6,T]$ use $\beta t\le c$. It follows that
\begin{equation}\label{eq:low-integral}
 \frac1{\pi\beta}\int_0^T\frac{|h(t)-e^{-t^2/2}|}{t}\,dt
       \le\frac{e^{25/54}}{6\pi}(I+J),
\end{equation}
where
\[
 I=\int_0^6t^2e^{-t^2/2+(269/5000)t^3}\,dt,
 \qquad J=\int_6^\infty t^2e^{-3t^2/25}\,dt.
\]

For the high frequencies, $|h_W(t)|$ is bounded by the nonnegative
characteristic function of a symmetric infinitely divisible variable
$\Pi$ with L\'evy measure $\frac12\sum_i\Law(X_i-X_i')$ off zero.
This follows from $|h_i|\le\exp[-(1-|h_i|^2)/2]$.
Its variance is $V=1-a$ and its third jump integral is at most $2\beta$.
By \eqref{eq:cp-bound},
$\Delta(\Pi,N(0,V))\le0.6102\beta/V^{3/2}$.
Integration by parts against a Gaussian density of variance $a/2$ gives
\begin{align*}
 f_{\Pi+G_{a/2}}(0)
 &\le\frac{\phz}{\sqrt{V+a/2}}
      +\frac{0.6102\beta}{V^{3/2}}
                          \frac{\sqrt{2/\pi}}{\sqrt{a/2}}\\
 &<\frac{\phz}{\sqrt{0.95}}
      +\frac{0.6102}{0.9^{3/2}}\frac{17}{20}\frac2{\sqrt\pi}
       <\frac{11}{10}.
\end{align*}
The first inequality uses the total variation of the Gaussian density;
the second uses \eqref{eq:gauss-window}. The nonnegative characteristic
function of $\Pi$ and exact density inversion now give
\[
 \frac1\pi\int_T^\infty\frac{e^{-at^2/2}|h_W(t)|}{t}\,dt
 \le\frac{e^{-aT^2/4}}{T} f_{\Pi+G_{a/2}}(0).
\]
Consequently its ratio to $\beta$ is less than
$(11/19)e^{-361/289}<1/6$. The corresponding normal tail is at most
\[
 \frac1{\pi\beta}\int_T^\infty\frac{e^{-t^2/2}}t\,dt
 \le\frac{\beta}{\pi c^2}e^{-c^2/(2\beta^2)}<10^{-9}.
\]

Appendix~\ref{app:rational} supplies the complete rational checks
\[
 I<23/10,\quad J<3/8,\quad e^{25/54}/(6\pi)<17/200,
\]
and the two scalar high-frequency comparisons just used. Exact CDF
inversion is applicable because the sum has Gaussian variance $a>0$.
Combining its two frequency ranges proves
\[
 \frac{\Delta(S)}\beta
 <\frac{17}{200}\left(\frac{23}{10}+\frac38\right)
       +\frac16+10^{-9}
 =\frac{1182125003}{3000000000}<\frac25.
\]
\end{proof}

Equations \eqref{eq:gauss-window} and
Lemma~\ref{lem:gauss-window-exclusion} contradict $C>\ce>2/5$.
This proves $a=0$ in Theorem~\ref{thm:structure}.

\section{An infinite maximizing sum is nonatomic}\label{sec:nonatomic}
We now have $S=\sum_iX_i$, with uniformly bounded supports whose radii
tend to zero. Suppose this full law has an atom. We prove that only
finitely many variables can be nonconstant.

Choose a modal state $c_i$ of probability $q_i$, and set $p_i=1-q_i$.
Lemma~\ref{lem:atoms-modal} gives $\sum_i p_i<\infty$. Centering gives
\[
 |c_i|\le\sqrt{v_i p_i/q_i},\qquad
 \sum_i|c_i|\le\sqrt{\sum_i v_i}\sqrt{\sum_i p_i/q_i}<\infty,
\]
since $q_i\ge1/3$. Write $X_i=c_i+Y_i$. With probability one only
finitely many $Y_i$ are nonzero, and every finite configuration of
positive-probability states is an atom of the full law, because
$\prod_iq_i>0$.

Two compactness observations retain the endpoint conventions needed
below. For fixed $M$, the set of sums using at most $M$ nonzero $Y_i$
is compact. In a sequence of such sums, each of its at most $M$ selected
indices can be made either constant or divergent; divergent indices
contribute zero in the limit. Every limit is another finite configuration.
If all the selected nonzero values are positive, then, for each $b>0$,
these sums avoid an interval $[b-\eta,b)$ for some $\eta>0$.
Otherwise a sequence approaching $b$ from below has fixed terms whose
sum is $b$ and escaping positive terms, so it is eventually at least
$b$, a contradiction.

If the maximizing threshold $z$ is itself an atom of mass $m>0$, all
sufficiently late gaps from the mode to another state are negative.
Indeed, for a positive gap $d_i=u_i-c_i\to0$, the modal-branch CDF
$H_i(w)=\Pp(S-X_i+c_i\le w)$ satisfies
$\|H_i-F\|_\infty\le p_i\to0$. Hence
\[
 H_i(z-d_i)-H_i(z)\longrightarrow-m.
\]
The stationarity equation \eqref{eq:transfer-moments} instead makes
this difference tend to zero, since both states and their gap tend to
zero. This excludes such a sequence.

Fix $K$ so that $q_j\ge3/4$ for $j\ge K$, and select a nonmodal
state of an arbitrarily late nonconstant coordinate $i\ge K$. Set
$d_i=u_i-c_i$ and $S_i^0=S-X_i+c_i$. Its exact mass-transfer derivative is
\begin{equation}\label{eq:pivotal}
 g_i=-\operatorname{sgn}(d_i)\Pp(E_i),\qquad
 E_i=\begin{cases}
       \{z-d_i<S_i^0\le z\},&d_i>0,\\
       \{z<S_i^0\le z-d_i\},&d_i<0.
      \end{cases}
\end{equation}
Let $N_i$ count the nonmodal outcomes among $j\ge K$, $j\ne i$.
For every fixed integer $M$,
\begin{equation}\label{eq:large-count}
                  E_i\subset\{N_i>M\}
                  \quad\text{for all sufficiently late }i.
\end{equation}
If $z$ is not an atom, it has positive distance from the compact set of
all configurations with at most $M+K$ nonmodal states. The intervals in
\eqref{eq:pivotal} shrink to $z$, proving \eqref{eq:large-count}.
If $z$ is an atom, take $K$ beyond the positive-gap exclusion. Conditional
on any of the finitely many early prefix values, write
$S_i^0=b-\sum_{j\ge K,j\ne i}Z_j$ with $Z_j\ge0$. Its pivotal interval
is exactly
\[
       \sum_{j\ge K,j\ne i}Z_j\in[b-z-|d_i|,b-z).
\]
For $b-z\le0$ it is empty; otherwise the one-sided compactness
observation proves \eqref{eq:large-count}. The right endpoint is
excluded, as required.

Vary every nonmodal probability for $j\ge K$, $j\ne i$, by the factor
$1+t$, taking the compensating mass from the mode. Both signs of
$|t|<1/2$ are admissible. The original product-law score is
\[
 L_i=\sum_{j\ge K,j\ne i}\ell_j,
 \qquad \ell_j=\begin{cases}1,&Y_j\ne0,\\-p_j/q_j,&Y_j=0.
 \end{cases}
 \qquad L_i\ge N_i-\Lambda,
 \quad\Lambda=\sum_{j\ge K}p_j/q_j<\infty.
\]
This infinite product variation is twice differentiable for every bounded
observable. Its density is
\[
       (1+t)^{N_i}\prod_{Y_j=0}(1-tp_j/q_j).
\]
The product and its first two derivatives are dominated, uniformly in
$i$, by a constant times a polynomial in $N_i$ times $(3/2)^{N_i}$.
These bounds are integrable because
$\E b^{N_i}\le\exp((b-1)\sum_{j\ge K}p_j)$ for every $b\ge1$.
Differentiating \eqref{eq:pivotal} in this direction gives the exact
mixed raw-CDF derivative
\[
 g_{i,t}=-\operatorname{sgn}(d_i)\E[\ind_{E_i}L_i].
\]
By \eqref{eq:large-count}, $\E[L_i\mid E_i]\to\infty$ whenever
these events have positive probability.

Let $h_i=\max\{|x|:x\in\supp X_i\}\to0$.
The mass-transfer stationarity equation and
\[
 |u_i^2-c_i^2|\le2h_i|d_i|,
 \qquad |r_i'|\le6h_i^2|d_i|
\]
give
\begin{equation}\label{eq:pivotal-size}
          g_i/d_i\to-\phi(z),\qquad
          \Pp(E_i)/|d_i|\to\phi(z)>0.
\end{equation}
It follows that $|g_{i,t}|/|d_i|\to\infty$.

On the other hand, form \eqref{eq:objective} with two parameters: the
selected mass transfer $s$ in coordinate $i$ and the collective variation
$t$ just defined. Uniformly in late $i$, its Hessian satisfies
\begin{equation}\label{eq:atomic-H}
       |J_{ss}|\le K_1d_i^2,\qquad
       |J_{st}-g_{i,t}|\le K_2|d_i|,\qquad
       |J_{tt}|\le K_3.
\end{equation}
To check these bounds, the selected derivatives are
$m_s=d_i$, $B_s=u_i^2-c_i^2$, $B_{ss}=-2d_i^2$,
$F_{ss}=0$, $F_s=O(|d_i|)$ and $|R_{ss}|\le18h_id_i^2$ by
\eqref{eq:transfer-moments}. All collective moment derivatives through
order two are bounded by summability of $p_j$ and the common support
bound. The raw CDF derivatives are uniformly bounded by the product
estimate above. Mixed moment derivatives vanish because the two
parameters affect disjoint coordinates. The normal normalization is
smooth near $(B,m)=(1,0)$, proving \eqref{eq:atomic-H}.
Both parameters are two-sided, so negative semidefiniteness gives
$|J_{st}|^2\le K_1K_3d_i^2$. This contradicts the divergent mixed CDF
ratio. Therefore an infinite maximizing sum is nonatomic.

Conversely, if $\sum_i(1-q_i)<\infty$, the same modal-drift argument
and Borel--Cantelli show that the series is a countable discrete law.
This proves the modal-deficit assertion and completes the proof of
Theorem~\ref{thm:structure}.

\section{Support motions and the completion rank}\label{sec:rank}
We prove Theorem~\ref{thm:rank} directly at the maximizing
representation. Suppose, to the contrary, that finite sets of
arbitrarily late coordinates admit nonconstant state functions with
constant sum on their completion sets. Subtract their individual means,
absorbing the sum into the constant $c$. Write
\[
 \E f_i(X_i)=0,\quad w_i=\E f_i(X_i)^2,
 \quad W=\sum_{i\in I}w_i>0,
 \quad V=\sum_{i\in I}v_i,
 \quad h=\max_{i\in I}\max|\supp X_i|.
\]
Along the supposed sequence, $h+V\to0$.

Move each raw state $x_i$ to $x_i+t f_i(x_i)$ and move the raw threshold
to $z+tc$. Every completion pattern has the same residual threshold as
before. For each other pattern the residual threshold lies outside the
closed support of $T_I$, in an open interval where its CDF is constant.
There are finitely many patterns, so one two-sided interval of sufficiently
small $t$ works for all of them. Distinct coordinate states remain distinct
after reducing the interval if necessary. Thus the raw CDF is exactly
unchanged, also when a completion is an atom or the endpoint of a gap.
This reasoning permits noncompact remainder support.

Choose $j$ maximizing $w_j/v_j$ among the nonconstant functions, and put
\[
 \begin{gathered}
 X=X_j,\quad f=f_j(X_j),\quad v=v_j,\quad w=w_j,\\
 m=\E Xf,\quad a_1=\E X^2f,\quad b_1=\E Xf^2,\quad
 M=\sum_{i\in I}\E[X_if_i(X_i)].
 \end{gathered}
\]
Then
\begin{equation}\label{eq:rank-selection}
                         v/w\le V/W.
\end{equation}
Tilt the original law of coordinate $j$ to $(1+sf)\,d\Law(X)$, and
apply the support motion. Both signs of $s$ are admissible because its
support is finite. Its moved mean is $sm+stw$. The centered threshold
$x$ and total variance $B$ have derivatives
\begin{align}
 x_s=-m,\quad x_t=c,\quad x_{st}=-w,&\qquad
 B_s=a_1,\quad B_{ss}=-2m^2,\notag\\
 B_t=2M,\quad B_{tt}=2W,&\qquad B_{st}=2b_1.           \label{eq:rank-moments}
\end{align}
The exact raw CDF satisfies $F_{ss}=F_t=F_{st}=F_{tt}=0$.
Cauchy--Schwarz gives
\[
 |m|\le\sqrt{vw},\quad |a_1|\le h\sqrt{vw},\quad
 |b_1|\le hw,\quad |M|\le\sqrt{VW}.
\]

For clarity, the separately recentered third moments have the following
derivatives; all expectations involving $X,f$ concern the selected
coordinate:
\begin{align}
 R_t&=3\sum_{i\in I}\E[X_i|X_i|f_i],&
 R_{tt}&=6\sum_{i\in I}\E[|X_i|f_i^2],\notag\\
 R_s&=\E[|X|^3f]-3m\E[X|X|],&
 R_{ss}&=6m^2\E|X|-6m\E[X|X|f],\notag\\
 R_{st}&=3\E[X|X|f^2]-3w\E[X|X|]
                         -6m\E[|X|f].              \label{eq:rank-cubic}
\end{align}
It follows that
\begin{align*}
 |R_t|&\le3h\sqrt{VW},& |R_{tt}|&\le6hW,\\
 |R_s|&\le4h^2\sqrt{vw},& |R_{ss}|&\le12hvw,
 & |R_{st}|&\le12h^2w.
\end{align*}
These bounds use $\E X^4\le h^2v$ and do not require a lower bound
on any state probability.

First stationarity in \eqref{eq:objective} gives
\[
       \phi c=2AM-CR_t,\qquad
       F_s=-\phi m-Aa_1+CR_s.
\]
Thus $|c|\le K\sqrt{VW}$ and $|F_s|\le K\sqrt{vw}$.
Using \eqref{eq:rank-moments}--\eqref{eq:rank-cubic} in its Hessian
now gives
\[
 |J_{ss}|\le K_1vw,\qquad |J_{tt}|\le K_2W,\qquad
 J_{st}=\phi w+O((h+V)w).
\]
Indeed the term $-\phi x_{st}$ is $\phi w$, direct variance and
third-moment terms are $O(hw)$, and products of first derivatives are
$O(\sqrt{VWvw})=O(Vw)$ by \eqref{eq:rank-selection}.
All constants depend only on the fixed maximizing representation and
threshold. Negative semidefiniteness would imply, for sufficiently late
sets,
\[
              \phi^2w^2/4\le K_1K_2vwW\le K_1K_2Vw^2,
\]
contradicting $V\to0$. This proves Theorem~\ref{thm:rank}.

\begin{proof}[Proof of Corollary~\ref{cor:separated}]
Fix a late cut and absorb all earlier blocks into the finite prefix;
let it have $N$ distinct sum values. Select $M>2N$ consecutive remaining
blocks. The unchanged infinite future has support in an interval of
positive width $W=\sum_{l\text{ after these blocks}}H_l$.
Distinct vectors of selected block-sum values have total sums at least
$W$ apart: at their first differing block $k$, their difference in
absolute value is at least
\[
       \delta_k-\sum_{l>k,\ l\text{ selected}}H_l\ge W.
\]
For each prefix atom, at most two such vectors put the maximizing
threshold inside the translated closed future interval. Thus at most
$2N$ selected block vectors can complete through the actual remainder
support, including endpoints.

There are $M+1$ unknowns $a_1,\ldots,a_M,c$ and at most $2N$
homogeneous equations $\sum_k a_k u_k=c$ on these vectors. A solution
exists with some $a_k\ne0$. Apply Theorem~\ref{thm:rank} to the
individual functions $f_i(x)=a_kx$ for coordinates in block $k$.
At least one is nonconstant because each block is nonconstant.
Their sum is constant on every completion pattern, a contradiction for
a sufficiently late cut. The total variance and largest coordinate
radius of the selected finite family tend to zero, independent of $M$.
\end{proof}

\appendix
\section{Rational checks for the Gaussian Fourier bound}\label{app:rational}
The finite calculation in Lemma~\ref{lem:gauss-window-exclusion} uses
only rational arithmetic. For $x\ge0$ define
\[
 E_n(x)=\sum_{k=0}^n\frac{x^k}{k!},\qquad
 U_n(x)=E_n(x)+\frac{x^{n+1}}{(n+1)!}
                         \frac1{1-x/(n+2)}\quad(0\le x<n+2).
\]
Then $E_n(x)\le e^x\le U_n(x)$, by the positive exponential series
and a geometric bound on its remaining terms.
The exponent $-t^2/2+(269/5000)t^3$ is decreasing on $[0,6]$.
With $d=3/100$, $l_j=jd$, $r_j=(j+1)d$, the exact upper sum
\begin{equation}\label{eq:finite-sum}
 \sum_{j=0}^{199}
 \frac{d\,r_j^2}{E_{40}(l_j^2/2-(269/5000)l_j^3)}<\frac{23}{10}
\end{equation}
therefore bounds $I$. Integration by parts and the normal Mills bound
give
\[
          J\le\frac{6025}{216}e^{-108/25}
             \le\frac{6025}{216E_{40}(108/25)}<\frac38.
\]
The remaining rational comparisons are
\begin{align*}
 \frac{U_{24}(25/54)}{6(157/50)}&<\frac{17}{200},&
 \frac{11}{19E_{40}(361/289)}&<\frac16,\\
 \frac{25}{61}+
 \frac{(3051/5000)(17/20)2}
      {(443/250)(9/10)(237/250)}&<\frac{11}{10}.&&
\end{align*}
The last expression bounds the Gaussian-convolution density because
\[
 (61/25)^2<(19/10)(157/50),\quad
 (443/250)^2<157/50,\quad (237/250)^2<9/10.
\]
The normal tail bound follows from $T>7$, $\beta<3/10$, $\pi>3$,
$c>1$, $e>5/2$, and $(5/2)^{24}>10^8$.
All these are finite rational inequalities, including
\eqref{eq:finite-sum}; the executable companion
\texttt{gaussian\_window\_certificate.py} evaluates them using Python's
standard-library rational type and asserts each comparison. It contains
no floating-point quadrature or external numerical library. Its final
exact comparison is
\[
          \frac{1182125003}{3000000000}<\frac25.
\]
The classical numerical constant $0.3051$ in
\eqref{eq:cp-bound} is an external theorem of Korolev and Shevtsova;
this finite check verifies the additional comparisons used here.

\bibliographystyle{plainnat}
\bibliography{refs}
\end{document}
